% --- Archivo: 12_restricciones_mixtas_penalizacion.tex ---

\subsection{Métodos de penalización cuadrática}\label{v2:sec:4.5.3}

De nuevo, consideramos el problema
\begin{equation}
    \min f(x) \quad \text{sujeto a} \quad x \in D = \left\{ x \in \R^n \;\middle|\; \begin{array}{l} h(x) = 0, \\ g(x) \le 0. \end{array} \right\},
    \label{v2:eq:4.181}
\end{equation}
donde $f \colon \R^n \to \R$, $h \colon \R^n \to \R^l$ y $g \colon \R^n \to \R^m$ son funciones diferenciables.

A seguir, extenderemos los resultados sobre la penalización cuadrática, obtenidos en la \cref{v2:sec:4.4.2} para restricciones de igualdad, al caso de restricciones mixtas del problema \eqref{v2:eq:4.181}. Para esto, utilizaremos la técnica de introducción de variables de holgura. Acontece que la penalización cuadrática, así como Lagrangianas aumentadas (que serán consideradas más adelante), son uno de los pocos casos donde la introducción de variables de holgura es una técnica eficiente, y no lleva a resultados más débiles de los que pueden ser obtenidos considerando desigualdades directamente.

Como ya sabemos, un problema con restricciones de igualdad y de desigualdad, como \eqref{v2:eq:4.181}, puede ser formalmente transformado en un problema con solamente restricciones de igualdad. Por ejemplo, de la siguiente manera:
\begin{equation}
    \min f(x) \quad \text{sujeto a} \quad (x, t) \in \tilde{D},
    \label{v2:eq:4.182}
\end{equation}
\begin{equation}
    \tilde{D} = \left\{ (x, t) \in \R^n \times \R^m \;\middle|\; \begin{array}{l} h(x) = 0, \\ g_i(x) + t_i^2 = 0, \; i = 1, \dots, m. \end{array} \right\}
    \label{v2:eq:4.183}
\end{equation}

Notamos, primero, que para todo $x \in \R^n$, podemos considerar que la condición $(x, t) \in \tilde{D}$ define el vector $t$ unívocamente, si no nos preocupamos con los signos de las coordenadas del vector $t$. Este punto de vista tiene sentido, porque la elección de estos signos no tiene ninguna influencia en el valor de la función objetivo del problema \eqref{v2:eq:4.182}-\eqref{v2:eq:4.183}, en el punto $(x, t)$. Por eso, a lo largo del análisis a seguir, vamos a ignorar los signos de las coordenadas de $t$. O sea, cualesquier dos vectores $(x, t^1)$ y $(x, t^2)$, que sean diferentes solamente en signos de las coordenadas de $t^1$ y $t^2 \in \R^m$, serán pensados como ``iguales'' (o, aún, como un mismo punto). Adoptando este punto de vista, soluciones locales de los problemas \eqref{v2:eq:4.181} y \eqref{v2:eq:4.182}-\eqref{v2:eq:4.183}, unívocamente corresponden una al otro: el punto $\bar{x} \in \R^n$ es una solución local de \eqref{v2:eq:4.181} si, y solamente si, el punto $(\bar{x}, \bar{t})$, donde
\begin{equation}
    \bar{t} = (\sqrt{-g_1(\bar{x})}, \dots, \sqrt{-g_m(\bar{x})}),
    \label{v2:eq:4.184}
\end{equation}
está bien definido, y es una solución local del problema \eqref{v2:eq:4.182}-\eqref{v2:eq:4.183} (vea el Ejercicio 4.5.7 más adelante). Más aún, el problema \eqref{v2:eq:4.182}-\eqref{v2:eq:4.183} no posee otras soluciones locales de la forma $(\bar{x}, t), t \in \R^m$.

Consideremos el término de penalización cuadrática para el problema \eqref{v2:eq:4.181}, i.e.,
\begin{equation}
    \psi(x) = \frac{1}{2} \left( \|h(x)\|^2 + \sum_{i=1}^m (\max\{0, g_i(x)\})^2 \right), \quad x \in \R^n,
    \label{v2:eq:4.185}
\end{equation}
donde introducimos el multiplicador $1/2$ por conveniencia en la hora de tomar derivadas. La familia asociada de funciones penalizadas viene dada por
\begin{equation}
    \varphi(\cdot; c) \colon \R^n \to \R, \quad \varphi(x; c) = f(x) + c\psi(x),
    \label{v2:eq:4.186}
\end{equation}
y los problemas penalizados asociados son
\begin{equation}
    \min \varphi(x; c), \quad x \in \R^n.
    \label{v2:eq:4.187}
\end{equation}
El método en cuestión, entonces, es el siguiente.

\noindent \textbf{Algoritmo 4.5.3. (El método de penalización cuadrática)} \\
Elegir $c_0 > 0$ y tomar $k := 0$.
\begin{enumerate}
    \item Calcular $x^k \in \R^n$, un punto estacionario del problema \eqref{v2:eq:4.187}, donde la función objetivo es dada por \eqref{v2:eq:4.186} con $c = c_k$.
    \item Elegir $c_{k+1} > c_k$. Tomar $k := k + 1$ y retornar al Paso 1.
\end{enumerate}

Los comentarios sobre la aplicación de este método son, en gran parte, los mismos que en el caso de la penalización cuadrática de restricciones de igualdad en la \cref{v2:sec:4.4.2}. Existe, sin embargo, una diferencia que merece ser mencionada. La función $\varphi(\cdot; c)$, definida por \eqref{v2:eq:4.186}-\eqref{v2:eq:4.185}, es diferenciable cuando $f, h$ y $g$ son diferenciables. No obstante, ella no es diferenciable dos veces en puntos $x \in \R^n$ para los cuales el conjunto de índices $I(x) = \{i = 1, \dots, m \mid g_i(x) = 0\}$ es no vacío, y esto es independiente de la diferenciabilidad de $f, h$ y $g$. Esto tiene que ser llevado en consideración en la hora de escoger un método de optimización irrestricta para la resolución de los problemas \eqref{v2:eq:4.187}.

Los Teoremas 4.5.4 y 4.5.5, que serán presentados a seguir, generalizan para el caso de restricciones mixtas los Teoremas 4.4.2 y 4.4.4, respectivamente, que fueron probados para restricciones de igualdad. Es más fácil conseguir estas generalizaciones apelando a la técnica de introducción de variables de holgura. Para eso, precisaremos de resultados auxiliares sobre la relación entre el problema \eqref{v2:eq:4.181} y el problema \eqref{v2:eq:4.182}-\eqref{v2:eq:4.183}.

% [Ejercicio 4.5.7 movido a exercises.tex]

Introducimos también la Lagrangiana del problema original \eqref{v2:eq:4.181} y la Lagrangiana para el problema modificado \eqref{v2:eq:4.182}-\eqref{v2:eq:4.183}:
\[
    \tilde{L} \colon (\R^n \times \R^m) \times (\R^l \times \R^m) \to \R,
\]
\begin{align*}
    \tilde{L}((x, t), (\lambda, \mu)) &= f(x) + \langle \lambda, h(x) \rangle + \sum_{i=1}^m \mu_i (g_i(x) + t_i^2) \\
    &= L(x, \lambda, \mu) + \sum_{i=1}^m \mu_i t_i^2.
\end{align*}

\begin{lemma}\label{v2:lem:4.5.1}
    Sean $f \colon \R^n \to \R$, $h \colon \R^n \to \R^l$ y $g \colon \R^n \to \R^m$ funciones diferenciables en el punto $\bar{x} \in \R^n$.
    \begin{enumerate}[(a)]
        \item Si $\bar{x}$ es un punto estacionario del problema \eqref{v2:eq:4.181} con multiplicadores de Lagrange asociados $\bar{\lambda} \in \R^l$ y $\bar{\mu} \in \R^m_+$, entonces el punto $(\bar{x}, \bar{t})$, donde $\bar{t} \in \R^m$ es dado por \eqref{v2:eq:4.184}, está bien definido y es estacionario para el problema \eqref{v2:eq:4.182}-\eqref{v2:eq:4.183}, con multiplicadores de Lagrange asociados $(\bar{\lambda}, \bar{\mu})$, i.e., $(\bar{x}, \bar{t}) \in \tilde{D}$ y
        \[
            \tilde{L}'_{(x,t)}((\bar{x}, \bar{t}), (\bar{\lambda}, \bar{\mu})) = 0.
        \]
        \item Si para algún $\bar{t} \in \R^m$ el punto $(\bar{x}, \bar{t})$ es estacionario para el problema \eqref{v2:eq:4.182}-\eqref{v2:eq:4.183}, con multiplicadores de Lagrange asociados $(\bar{\lambda}, \bar{\mu}) \in \R^l \times \R^m$, tem-se que $\bar{x} \in D$ y
        \begin{equation}
            L'_x(\bar{x}, \bar{\lambda}, \bar{\mu}) = 0,
            \label{v2:eq:4.188}
        \end{equation}
        \begin{equation}
            \bar{\mu}_i g_i(\bar{x}) = 0, \quad i = 1, \dots, m.
            \label{v2:eq:4.189}
        \end{equation}
    \end{enumerate}
\end{lemma}
\begin{proof}
    El hecho de que $\bar{x}$ es un punto estacionario del problema \eqref{v2:eq:4.181}, con multiplicadores de Lagrange $\bar{\lambda} \in \R^l$ y $\bar{\mu} \in \R^m_+$, significa la validez de las restricciones, de la igualdad \eqref{v2:eq:4.188}, y de las condiciones de complementaridad \eqref{v2:eq:4.189}. De \eqref{v2:eq:4.184} y \eqref{v2:eq:4.190} sigue que el punto $(\bar{x}, \bar{t})$ es un punto factible del problema \eqref{v2:eq:4.182}-\eqref{v2:eq:4.183}. Además de eso, de \eqref{v2:eq:4.188} y \eqref{v2:eq:4.189} obtenemos que
    \begin{align*}
        \tilde{L}'_{(x,t)}((\bar{x}, \bar{t}), (\bar{\lambda}, \bar{\mu})) &= L'_x(\bar{x}, \bar{\lambda}, \bar{\mu}) = 0, \\
        \tilde{L}'_{t_i}((\bar{x}, \bar{t}), (\bar{\lambda}, \bar{\mu})) &= 2\bar{\mu}_i \bar{t}_i = 0 \quad \forall i = 1, \dots, m.
    \end{align*}
    Esto concluye la prueba del ítem (a).
    El ítem (b) se prueba repitiendo el mismo raciocinio, en el orden inverso.
\end{proof}

Cabe resaltar que, en el ítem (b) del \cref{v2:lem:4.5.1}, no se afirma que $\bar{x}$ es un punto estacionario del problema \eqref{v2:eq:4.181} en el sentido completo: el multiplicador $\bar{\mu}$ puede poseer coordenadas negativas. Para un punto $x \in \R^n$ cualquiera, definimos el conjunto de índices de restricciones de desigualdad del problema \eqref{v2:eq:4.181} violadas en $x$, i.e.,
\[
    I^+(x) = \{i = 1, \dots, m \mid g_i(x) > 0\}.
\]
Decimos que $x$ satisface la condición de independencia lineal de los gradientes extendida, si
\begin{equation}
    \{h'_i(x), \; i = 1, \dots, l\} \cup \{g'_i(x), \; i \in I(x) \cup I^+(x)\}
    \label{v2:eq:4.191}
\end{equation}
es un conjunto linealmente independiente.

Notamos que esta definición el punto $x$ puede no ser factible. Cuando $x$ es factible, tem-se que $I^+(x) = \emptyset$, y la condición \eqref{v2:eq:4.191} introducida arriba coincide con la condición de regularidad clásica \eqref{v1:eq:1.17}.

% [Ejercicio 4.5.8 movido a exercises.tex]

Continuamos con el estudio de conexiones entre los problemas \eqref{v2:eq:4.181} y \eqref{v2:eq:4.182}-\eqref{v2:eq:4.183}.

\begin{lemma}\label{v2:lem:4.5.2}
    Sean $f \colon \R^n \to \R$, $h \colon \R^n \to \R^l$ y $g \colon \R^n \to \R^m$ funciones dos veces diferenciables en el punto $\bar{x} \in \R^n$. Supongamos que $\bar{x}$ satisfaga la condición de independencia lineal de los gradientes extendida \eqref{v2:eq:4.191}, y que $\bar{x}$ sea un punto estacionario del problema \eqref{v2:eq:4.181}, con multiplicadores de Lagrange asociados $\bar{\lambda} \in \R^l$ y $\bar{\mu} \in \R^m_+$. Supongamos, finalmente, que $\bar{x}$ satisfaga la condición suficiente de segunda orden del Teorema 1.2.2(b) y la condición de complementaridad estricta (1.25).
    Entonces el punto $(\bar{x}, \bar{t})$, donde $\bar{t} \in \R^m$ es dado por \eqref{v2:eq:4.184}, está bien definido y satisface la condición de regularidad de las restricciones para el problema \eqref{v2:eq:4.182}-\eqref{v2:eq:4.183}. Además de eso, $(\bar{x}, \bar{t})$ es un punto estacionario de este problema, con un único par de multiplicadores de Lagrange asociados $(\bar{\lambda}, \bar{\mu})$. Más aún, este punto satisface la condición suficiente de segunda orden para el problema \eqref{v2:eq:4.182}-\eqref{v2:eq:4.183}, i.e.,
    \begin{equation}
        \langle \tilde{L}''_{(x,t)(x,t)}((\bar{x}, \bar{t}), (\bar{\lambda}, \bar{\mu}))(d, \tau), (d, \tau) \rangle > 0
        \label{v2:eq:4.192}
    \end{equation}
    para todo $(d, \tau) \in (\R^n \times \R^m) \setminus \{0\}$ tal que
    \begin{equation}
        h'(\bar{x})d = 0, \quad \langle g'_i(\bar{x}), d \rangle + 2\bar{t}_i \tau_i = 0 \quad \forall i = 1, \dots, m.
        \label{v2:eq:4.193}
    \end{equation}
\end{lemma}
\begin{proof}
    El hecho de que $(\bar{x}, \bar{\sigma})$ es un punto estacionario con multiplicadores de Lagrange asociados $(\bar{\lambda}, \bar{\mu})$, fue establecido en el ítem (a) del \cref{v2:lem:4.5.1}. La validez en este punto de la condición de regularidad de las restricciones sigue del resultado del Ejercicio 4.5.8.
    Tomamos $(d, \tau) \in (\R^n \times \R^m) \setminus \{0\}$ que satisface \eqref{v2:eq:4.193}. Haciendo el cálculo directo, puede ser visto que
    \begin{align}
        &\langle \tilde{L}''_{(x,t)(x,t)}((\bar{x}, \bar{t}), (\bar{\lambda}, \bar{\mu}))(d, \tau), (d, \tau) \rangle \nonumber \\
        &\quad = \langle L''_{xx}(\bar{x}, \bar{\lambda}, \bar{\mu})d, d \rangle + 2\sum_{i=1}^m \bar{\mu}_i \tau_i^2.
        \label{v2:eq:4.194}
    \end{align}
    De \eqref{v2:eq:4.184}, \eqref{v2:eq:4.193}, y de la condición de complementaridad estricta (1.25), tenemos que $d \in \mathcal{K}(\bar{x})$, donde, como siempre,
    \[
        \mathcal{K}(\bar{x}) = \{d \in \ker h'(\bar{x}) \mid \langle g'_i(\bar{x}), d \rangle \le 0 \quad \forall i \in I(\bar{x}), \; \langle f'(\bar{x}), d \rangle \le 0\}
    \]
    es el cono crítico del problema \eqref{v2:eq:4.181} en el punto $\bar{x}$. Si $d \neq 0$, la desigualdad \eqref{v2:eq:4.192} sigue inmediatamente de la condición suficiente de segunda orden para el problema \eqref{v2:eq:4.181} en el punto $\bar{x}$ y de la igualdad \eqref{v2:eq:4.194}, donde el segundo término en el lado derecho es no-negativo (ya que $\bar{\mu} \ge 0$).
    Sea $d = 0$. En este caso, $\tau \neq 0$. De \eqref{v2:eq:4.184} y \eqref{v2:eq:4.193}, sigue que $\tau_i = 0 \; \forall i \in \{1, \dots, m\} \setminus I(\bar{x})$. Por tanto, entre los números $\tau_i, i \in I(\bar{x})$, hay algunos estrictamente positivos. De aquí, de la condición de complementaridad estricta (1.25), y de \eqref{v2:eq:4.194}, obtenemos la desigualdad \eqref{v2:eq:4.192}.
\end{proof}

Ahora, para el problema \eqref{v2:eq:4.182}-\eqref{v2:eq:4.183}, introducimos la familia de funciones penalizadas
\begin{equation}
    \tilde{\varphi}(\cdot; c) \colon \R^n \times \R^m \to \R, \quad \tilde{\varphi}(x, t; c) = f(x) + c\tilde{\psi}(x, t),
    \label{v2:eq:4.195}
\end{equation}
asociada a la penalización cuadrática
\begin{equation}
    \tilde{\psi} \colon \R^n \times \R^m \to \R, \quad \tilde{\psi}(x, t) = \frac{1}{2} \left( \|h(x)\|^2 + \sum_{i=1}^m (g_i(x) + t_i^2)^2 \right).
    \label{v2:eq:4.196}
\end{equation}
Consideraremos los problemas penalizados correspondientes:
\begin{equation}
    \min \tilde{\varphi}(x, t; c), \quad (x, t) \in \R^n \times \R^m.
    \label{v2:eq:4.197}
\end{equation}
Observemos que para todo $x \in \R^n$ fijo, la minimización en relación a $t \in \R^m$ en \eqref{v2:eq:4.197} puede ser hecha explícitamente, eliminando así las variables auxiliares: el mínimo global se alcanza en $t(x) \in \R^m$, donde
\begin{equation}
    t_i(x) = \begin{cases}
        0, & \text{si } g_i(x) > 0, \\
        \sqrt{-g_i(x)}, & \text{caso contrario,}
    \end{cases} \quad i = 1, \dots, m.
    \label{v2:eq:4.198}
\end{equation}
Este minimizador es único, si ignoramos los signos de las coordenadas. Tenemos, por tanto, lo siguiente:
\begin{align*}
    \min_{t \in \R^m} \tilde{\varphi}(x, t; c) &= \tilde{\varphi}(x, t(x); c) \\
    &= f(x) + \frac{c}{2} \left( \|h(x)\|^2 + \sum_{i=1}^m (\max\{0, g_i(x)\})^2 \right) \\
    &= \varphi(x; c) \quad \forall x \in \R^n,
\end{align*}
donde $\varphi(\cdot; c)$ es la función definida por \eqref{v2:eq:4.186}, \eqref{v2:eq:4.185}.

% [Ejercicio 4.5.9 movido a exercises.tex]

\begin{lemma}\label{v2:lem:4.5.3}
    Sean $f \colon \R^n \to \R$, $h \colon \R^n \to \R^l$ y $g \colon \R^n \to \R^m$ funciones diferenciables en el punto $\bar{x} \in \R^n$.
    Entonces para todo $c \ge 0$, un punto $\bar{x}$ es estacionario para el problema \eqref{v2:eq:4.187} con función objetivo dada por \eqref{v2:eq:4.186}, \eqref{v2:eq:4.185}, si, y solamente si, el punto $(\bar{x}, t(\bar{x}))$, donde $t(\bar{x}) \in \R^m$ está dado por \eqref{v2:eq:4.198}, es estacionario para el problema \eqref{v2:eq:4.197} con función objetivo definida por \eqref{v2:eq:4.195}, \eqref{v2:eq:4.196}.
\end{lemma}
\begin{proof}
    Sea $\tilde{x}$ un punto estacionario del problema \eqref{v2:eq:4.187}. De \eqref{v2:eq:4.198} obtenemos que
    \begin{align*}
        \tilde{\varphi}'_x(\tilde{x}, t(\tilde{x}); c) &= f'(\tilde{x}) + c \left( (h'(\tilde{x}))^\top h(\tilde{x}) + \sum_{i=1}^m (g_i(\tilde{x}) + (t_i(\tilde{x}))^2) g'_i(\tilde{x}) \right) \\
        &= f'(\tilde{x}) + c \left( (h'(\tilde{x}))^\top h(\tilde{x}) + \sum_{i=1}^m \max\{0, g_i(\tilde{x})\} g'_i(\tilde{x}) \right) \\
        &= \varphi'(\tilde{x}; c) = 0,
    \end{align*}
    donde la última igualdad sigue de \eqref{v2:eq:4.186} y \eqref{v2:eq:4.185}. Por otro lado, de \eqref{v2:eq:4.198} tenemos
    \begin{equation}
        \tilde{\varphi}'_{t_i}(\tilde{x}, t(\tilde{x}); c) = 2c(g_i(\tilde{x}) + (t_i(\tilde{x}))^2) t_i(\tilde{x}) = 0 \quad \forall i = 1, \dots, m.
        \label{v2:eq:4.200}
    \end{equation}
    Las igualdades \eqref{v2:eq:4.199} y \eqref{v2:eq:4.200} significan que $(\tilde{x}, t(\tilde{x}))$ es un punto estacionario del problema \eqref{v2:eq:4.197}.
    La afirmación recíproca se prueba invirtiendo el raciocinio utilizado para obtener \eqref{v2:eq:4.199}.
\end{proof}

Notemos que el problema \eqref{v2:eq:4.197} puede poseer puntos estacionarios $(\tilde{x}, \tilde{t})$ tales que ni todos los valores absolutos de las coordenadas de $\tilde{t} \in \R^m$ coincidan con las coordenadas correspondientes de $t(\tilde{x})$, y en este caso $\tilde{x} \in \R^n$ puede no ser un punto estacionario del problema \eqref{v2:eq:4.187}.

\begin{example}\label{v2:exa:4.5.1}
    Sean $n = m = 1$, $l = 0$,
    \[
        f(x) = x, \quad g(x) = x^2 - 1, \quad x \in \R.
    \]
    Puntos estacionarios de la función $\tilde{\varphi}(\cdot; c)$ se caracterizan por las ecuaciones
    \begin{align*}
        \tilde{\varphi}'_x(x, t; c) &= 1 + 2c(x^2 + t^2 - 1)x = 0, \\
        \tilde{\varphi}'_t(x, t; c) &= 2c(x^2 + t^2 - 1)t = 0.
    \end{align*}
    Es fácil ver que, para todo $c$ suficientemente grande, este sistema tiene una solución de la forma $(\tilde{x}, 0)$, donde $\tilde{x} \in [-1, 1]$.
    Sin embargo,
    \[
        \varphi'(\tilde{x}; c) = 1 + 2c \max\{0, \tilde{x}^2 - 1\}\tilde{x} = 1,
    \]
    i.e., $\tilde{x}$ no es un punto estacionario de la función $\varphi_c$.
    Generalmente, el fenómeno expuesto arriba no presenta una dificultad grave, porque puntos estacionarios ``indeseados'' del problema \eqref{v2:eq:4.197} no pueden ser soluciones locales (o sea, minimizadores locales) de este problema.
\end{example}

\begin{lemma}\label{v2:lem:4.5.4}
    Sean $f \colon \R^n \to \R$, $h \colon \R^n \to \R^l$ y $g \colon \R^n \to \R^m$ cualesquiera. Si para algún $c > 0$ el punto $(\tilde{x}, \tilde{t}) \in \R^n \times \R^m$ es una solución local del problema \eqref{v2:eq:4.197} con función objetivo dada por \eqref{v2:eq:4.195}, \eqref{v2:eq:4.196}, entonces $\tilde{t}$ coincide con el punto $t(\tilde{x})$, definido de acuerdo con \eqref{v2:eq:4.198}, ignorando los signos de las coordenadas.
\end{lemma}
\begin{proof}
    Definimos la función
    \[
        \hat{\varphi}(\cdot; c) \colon \R^m \to \R,
    \]
    \[
        \hat{\varphi}(t; c) = \tilde{\varphi}(\tilde{x}, t; c) = f(\tilde{x}) + \frac{c}{2} \left( \|h(\tilde{x})\|^2 + \sum_{i=1}^m (g_i(\tilde{x}) + t_i^2)^2 \right).
    \]
    Esta función es (infinitamente) diferenciable en $\R^m$ y el punto $\tilde{t}$ es su minimizador irrestricto. Por el Teorema 1.2.1(a), tenemos que
    \begin{equation}
        \hat{\varphi}'_{t_i}(\tilde{t}; c) = 2c(g_i(\tilde{x}) + \tilde{t}_i^2)\tilde{t}_i = 0 \quad \forall i = 1, \dots, m.
        \label{v2:eq:4.201}
    \end{equation}
    De aquí sigue que para los índices $i = 1, \dots, m$, para los cuales $g_i(\tilde{x}) \ge 0$, tiene que ser $\tilde{t}_i = 0 = t_i(\tilde{x})$.
    Ahora, por el Teorema 1.2.1(a), la matriz $\hat{\varphi}''(\tilde{t}; c)$ es semidefinida positiva. Si admitimos que para algún $i = 1, \dots, m$ tal que $g_i(\tilde{x}) < 0$, tenemos que $\tilde{t}_i^2 \neq (t_i(\tilde{x}))^2 = -g_i(\tilde{x})$, de \eqref{v2:eq:4.201} sigue que $\tilde{t}_i = 0$. En este caso,
    \[
        \hat{\varphi}''_{t_i t_i}(\tilde{t}; c) = 2c(g_i(\tilde{x}) + 3\tilde{t}_i^2) = 2cg_i(\tilde{x}) < 0,
    \]
    en contradicción con el hecho de que la matriz $\hat{\varphi}''(\tilde{t}; c)$ es semidefinida positiva (notemos que la matriz en cuestión es diagonal).
\end{proof}

En particular, y a diferencia del caso de los puntos estacionarios, las soluciones locales de los problemas \eqref{v2:eq:4.187} y \eqref{v2:eq:4.197} están en correspondencia única (en el mismo sentido que en el caso de soluciones locales de los problemas \eqref{v2:eq:4.181} y \eqref{v2:eq:4.182}-\eqref{v2:eq:4.183}).

\begin{theorem}\label{v2:thm:4.5.4}
    Sean funciones $f \colon \R^n \to \R$, $h \colon \R^n \to \R^l$ y $g \colon \R^n \to \R^m$ diferenciables en $\R^n$, con derivadas continuas.
    Entonces, si la sucesión $\{x^k\}$ es generada por el Algoritmo 4.5.3, donde $\{c_k\} \to +\infty$ $(k \to \infty)$ y se utiliza la función $\psi \colon \R^n \to \R$ dada por \eqref{v2:eq:4.185}, todo punto de acumulación $\bar{x} \in \R^n$ de la sucesión $\{x^k\}$ que satisface la condición de independencia lineal de los gradientes extendida \eqref{v2:eq:4.191}, es un punto estacionario del problema \eqref{v2:eq:4.181}.
    Además de eso, para toda subsucesión $\{x^{k_j}\}$ que converge a $\bar{x}$ se tiene
    \begin{equation}
        \{c_{k_j} h(x^{k_j})\} \to \bar{\lambda} \quad (j \to \infty),
        \label{v2:eq:4.202}
    \end{equation}
    \begin{equation}
        c_{k_j} \max\{0, g_i(x^{k_j})\} \to \bar{\mu}_i \quad (j \to \infty) \quad \forall i = 1, \dots, m,
        \label{v2:eq:4.203}
    \end{equation}
    donde $\bar{\lambda} \in \R^l$ y $\bar{\mu} \in \R^m_+$ son los únicos multiplicadores de Lagrange asociados a $\bar{x}$.
\end{theorem}
\begin{proof}
    Consideremos una sucesión $\{(x^k, t(x^k))\}$, donde para todo $k$ el punto $t(x^k) \in \R^m$ está definido de acuerdo con \eqref{v2:eq:4.198}. Pelo Lema 4.5.3, la sucesión $\{(x^k, t(x^k))\}$ puede ser pensada como generada por el método de penalización cuadrática aplicado al problema \eqref{v2:eq:4.182}-\eqref{v2:eq:4.183}. Pela continuidad de la función $t(\cdot)$, sigue que $(\bar{x}, t(\bar{x}))$ es un punto de acumulación de esta sucesión, y el Ejercicio 4.5.8 sigue que las restricciones del problema \eqref{v2:eq:4.182}-\eqref{v2:eq:4.183} son regulares en este punto.
    Utilizando el Teorema 4.4.2, obtenemos que $(\bar{x}, t(\bar{x}))$ es un punto estacionario del problema \eqref{v2:eq:4.182}-\eqref{v2:eq:4.183}, y que para una subsucesión $\{x^{k_j}\}$ convergente a $\bar{x}$ valen las relaciones \eqref{v2:eq:4.202} y
    \begin{equation}
        c_{k_j}(g_i(x^{k_j}) + (t_i(x^{k_j}))^2) \to \bar{\mu}_i \quad \forall i = 1, \dots, m \quad (j \to \infty),
        \label{v2:eq:4.204}
    \end{equation}
    donde $(\bar{\lambda}, \bar{\mu}) \in \R^l \times \R^m$ son los multiplicadores de Lagrange correspondientes. De \eqref{v2:eq:4.198}, sigue que \eqref{v2:eq:4.204} coincide con \eqref{v2:eq:4.203}, y de \eqref{v2:eq:4.203} sigue que $\bar{\mu} \ge 0$.
    La demostración se concluye llevando en cuenta el ítem (b) del Lema 4.5.1.
\end{proof}

Notemos que de \eqref{v2:eq:4.203} sigue la siguiente propiedad interesante: para todo $j$ suficientemente grande y todo $i \in I(\bar{x})$ tal que $\bar{\mu}_i > 0$, se tiene que $g_i(x^{k_j}) > 0$. Este hecho es consistente con las propiedades generales de los métodos de penalización externa (vea el § 4.3.1), de acuerdo con las cuales las iteraciones no son factibles en relación al problema original.
El Teorema 4.5.4 nos dice que todo punto de acumulación de una sucesión generada por el método de penalización cuadrática o es un punto estacionario del problema original \eqref{v2:eq:4.181}, o viola la condición de independencia lineal de los gradientes.
En este último caso, el punto en cuestión puede no ser factible para
el problema.

\begin{example}\label{v2:exa:4.5.2}
    Sean $n = m = 1, l = 0$,
    \[
        f(x) = x, \quad g(x) = x^2 + 1, \quad x \in \R.
    \]
    Tenemos que
    \[
        \varphi(x; c) = x + c(\max\{0, x^2 + 1\})^2 / 2 = x + c(x^2 + 1)^2 / 2, \quad x \in \R,
    \]
    y los puntos estacionarios del problema \eqref{v2:eq:4.187} se caracterizan por la ecuación
    \[
        \varphi'(x; c) = 1 + 2c(x^2 + 1)x = 0.
    \]
    Para todo $c$, esta ecuación tiene una única solución, que tiende a cero cuando $c \to \infty$. Mas el punto $0$ no es factible en el problema \eqref{v2:eq:4.181}. Observemos que $g'(0) = 0$, i.e., la condición \eqref{v2:eq:4.191} está violada.

    En el ejemplo anterior el problema original tiene su conjunto factible vacío. Más la violación de restricciones por un punto de acumulación de una sucesión generada por el método de penalización cuadrática también es posible en el caso de que el conjunto factible es no vacío (naturalmente, tal punto debe violar la condición \eqref{v2:eq:4.191} de independencia lineal de los gradientes).
\end{example}

\begin{example}\label{v2:exa:4.5.3}
    Sean $n = 2, l = m = 1$,
    \[
        h(x) = x_1^2 - 1, \quad g(x) = -x_1 + x_2^2, \quad x \in \R^2,
    \]
    y sea $f$ una función constante en $\R^2$.
    Puede ser ver que el punto $0$ es estacionario para el problema \eqref{v2:eq:4.187} para cualquier $c$, y que $0$ no es factible en el problema \eqref{v2:eq:4.181}. La condición \eqref{v2:eq:4.191} está violada, ya que $h'(0) = 0$.
\end{example}

Finalmente, presentamos estimativas de tasa de convergencia de los métodos de penalización cuadrática.

\begin{theorem}[Estimativas de tasa de convergencia de los métodos de penalización cuadrática]\label{v2:thm:4.5.5}
    Sean $f \colon \R^n \to \R, h \colon \R^n \to \R^l$ y $g \colon \R^n \to \R^m$ funciones dos veces diferenciables en una vecindad del punto $\bar{x} \in \R^n$, con derivadas segundas continuas en este punto. Supongamos que $\bar{x}$ satisface la condición \eqref{v2:eq:4.191} de independencia lineal de los gradientes, y que $\bar{x}$ sea un punto estacionario del problema \eqref{v2:eq:4.181}, con (únicos) multiplicadores de Lagrange $\bar{\lambda} \in \R^l$ y $\bar{\mu} \in \R^m_+$ asociados. Supongamos, finalmente, que $\bar{x}$ satisface la condición suficiente de segunda orden del Teorema 1.2.3(b) y la condición de complementaridad estricta (1.25).
    Entonces, existe una vecindad $U$ del punto $\bar{x}$ y números $\bar{c} \ge 0$ y $M > 0$, tales que, para todo $c > \bar{c}$, el problema \eqref{v2:eq:4.187} con función objetivo definida por las fórmulas \eqref{v2:eq:4.186}, \eqref{v2:eq:4.185}, posee en $U$ un único punto estacionario $x(c)$, y se tiene que
    \begin{equation}
        \|x(c) - \bar{x}\| \le M\frac{\|\bar{\lambda}\| + \|\bar{\mu}\|}{c}, \quad \|c f(x(c)) - \bar{\lambda}\| \le M\frac{\|\bar{\lambda}\| + \|\bar{\mu}\|}{c},
    \end{equation}
    \begin{equation}
        |c\max\{0, g_i(x(c))\} - \bar{\mu}_i| \le M\frac{\|\bar{\lambda}\| + \|\bar{\mu}\|}{c} \quad \forall i = 1, \dots, m.
        \label{v2:eq:4.205}
    \end{equation}
\end{theorem}

% [Ejercicios 4.5.10 y 4.5.11 movidos a exercises.tex]

De la condición de complementaridad estricta (1.25) y de la estimativa \eqref{v2:eq:4.205}, puede ser visto que para todo valor del parámetro penalizador $c$ suficientemente grande, todas las restricciones de desigualdad del problema \eqref{v2:eq:4.181}, activas en el punto $\bar{x}$, son violadas en los puntos de la trayectoria $x(\cdot)$ definida en el Teorema 4.5.5. O sea, tenemos que $g_i(x(c)) > 0 \; \forall i \in I(\bar{x})$ para todo $c$ suficientemente grande (comparar con las observaciones hechas para el Teorema 4.5.4).

\subsection{Lagrangianas aumentadas}\label{v2:sec:4.5.4}

A seguir, extenderemos los resultados sobre Lagrangianas aumentadas, obtenidos en la \cref{v2:sec:4.4.3} para restricciones de igualdad, al caso de restricciones mixtas del problema \eqref{v2:eq:4.181}. Para esto, de nuevo utilizaremos la técnica de introducción de variables de holgura.
Para todo $c > 0$, definimos la Lagrangiana aumentada para el problema \eqref{v2:eq:4.182}-\eqref{v2:eq:4.183} con restricciones transformadas en igualdades:
\[
    \tilde{L}(\cdot; c) \colon (\R^n \times \R^m) \times (\R^l \times \R^m) \to \R,
\]
\begin{align}
    &\tilde{L}((x, t), (\lambda, \mu); c) \nonumber \\
    &= \tilde{L}((x, t), (\lambda, \mu)) \nonumber \\
    &\quad + \frac{c}{2} \left( \|h(x)\|^2 + \sum_{i=1}^m (g_i(x) + t_i^2)^2 \right).
    \label{v2:eq:4.206}
\end{align}
Para $\lambda \in \R^l$ y $\mu \in \R^m$ dados, la minimización en
\begin{equation}
    \min \tilde{L}((x, t), (\lambda, \mu); c), \quad (x, t) \in \R^n \times \R^m,
    \label{v2:eq:4.207}
\end{equation}
en relación a $t \in \R^m$ puede ser hecha explícitamente. En particular, para todo $x \in \R^n$ fijo, el mínimo global se alcanza en el punto $t(x, \mu, c)$, donde
\begin{equation}
    t_i(x, \mu, c) = \begin{cases}
        0, & \text{se } g_i(x) + \mu_i/c > 0, \\
        \sqrt{-(g_i(x) + \mu_i/c)}, & \text{caso contrario},
    \end{cases}
    \label{v2:eq:4.208}
\end{equation}
$i = 1, \dots, m$. Este minimizador es único (ignorando los signos de las coordenadas).
Notemos que
\begin{align}
    g_i(x) + (t_i(x, \mu, c))^2 &= g_i(x) + \max\{0, -g_i(x) - \mu_i/c\} \nonumber \\
    &= \max\{g_i(x), -\mu_i/c\}.
    \label{v2:eq:4.209}
\end{align}
Donde, después de algunas transformaciones simples, obtenemos la Lagrangiana aumentada para el problema \eqref{v2:eq:4.181} de la siguiente forma:
\[
    L(\cdot; c) \colon \R^n \times \R^l \times \R^m \to \R,
\]
\begin{align}
    &L(x, \lambda, \mu; c) \nonumber \\
    &= \min_{t \in \R^m} \tilde{L}((x, t), (\lambda, \mu); c) \nonumber \\
    &= \tilde{L}((x, t(x, \mu, c)), (\lambda, \mu); c) \nonumber \\
    &= f(x) + \langle \lambda, h(x) \rangle + \frac{c}{2}\|h(x)\|^2 \nonumber \\
    &\quad + \frac{1}{2c} \sum_{i=1}^m \left( (\max\{0, c g_i(x) + \mu_i\})^2 - \mu_i^2 \right).
    \label{v2:eq:4.210}
\end{align}
Por tanto, en vez de \eqref{v2:eq:4.207}, consideramos el problema
\begin{equation}
    \min L(x, \lambda, \mu; c), \quad x \in \R^n.
    \label{v2:eq:4.211}
\end{equation}
El Teorema 4.5.6 a continuación muestra que puntos KKT del problema original \eqref{v2:eq:4.181} corresponden a puntos estacionarios irrestrictos de la Lagrangiana aumentada de este problema, para $c > 0$ cualquiera. Más primero precisaremos de un resultado auxiliar.

\begin{lemma}\label{v2:lem:4.5.5}
    Sean $a \in \R$ y $b \in \R$. Entonces
    \begin{equation}
        a = \max\{0, b\}
        \label{v2:eq:4.212}
    \end{equation}
    si, y solamente si,
    \begin{equation}
        a - b \ge 0, \quad a(a - b) = 0, \quad a \ge 0.
        \label{v2:eq:4.213}
    \end{equation}
\end{lemma}
\begin{proof}
    Observamos que \eqref{v2:eq:4.212} significa que $a \ge 0$, $a \ge b$, o $a = 0$ o $a = b$. Estas condiciones también pueden ser escritas como \eqref{v2:eq:4.213}.
\end{proof}

\begin{theorem}[Puntos estacionarios de la Lagrangiana aumentada]\label{v2:thm:4.5.6}
    Sean $f \colon \R^n \to \R$, $h \colon \R^n \to \R^l$ y $g \colon \R^n \to \R^m$ funciones diferenciables en el punto $\bar{x} \in \R^n$.
    Entonces para $\bar{\lambda} \in \R^l$ y $\bar{\mu} \in \R^m_+$, y $c > 0$ cualquier, el punto $(\bar{x}, \bar{\lambda}, \bar{\mu})$ es una solución del sistema de Karush-Kuhn-Tucker del problema \eqref{v2:eq:4.181} si, y solamente si,
    \[
        L'_x(\bar{x}, \bar{\lambda}, \bar{\mu}; c) = 0.
    \]
\end{theorem}
\begin{proof}
    Tenemos que
    \begin{align*}
        L'_x(\bar{x}, \bar{\lambda}, \bar{\mu}; c) &= f'(\bar{x}) + (h'(\bar{x}))^\top(\bar{\lambda} + c h(\bar{x})) \\
        &\quad + \sum_{i=1}^m (\max\{0, c g_i(\bar{x}) + \bar{\mu}_i\}) g'_i(\bar{x}),
    \end{align*}
    \[
        L'_\lambda(\bar{x}, \bar{\lambda}, \bar{\mu}; c) = h(\bar{x}),
    \]
    \[
        L'_{\mu_i}(\bar{x}, \bar{\lambda}, \bar{\mu}; c) = \frac{1}{c} (\max\{0, c g_i(\bar{x}) + \bar{\mu}_i\} - \bar{\mu}_i), \quad i = 1, \dots, m.
    \]
    En particular, $L'_\lambda(\bar{x}, \bar{\lambda}, \bar{\mu}; c) = 0$ si, y solamente si, $h(\bar{x}) = 0$. Como $c > 0$, tenemos también que $L'_{\mu_i}(\bar{x}, \bar{\lambda}, \bar{\mu}; c) = 0$ si, y solamente si,
    \[
        \bar{\mu}_i = \max\{0, c g_i(\bar{x}) + \bar{\mu}_i\}, \quad i = 1, \dots, m.
    \]
    Lo que, por el \cref{v2:lem:4.5.5}, es equivalente a
    \[
        -cg_i(\bar{x}) \ge 0, \quad \bar{\mu}_i(-cg_i(\bar{x})) = 0, \quad \bar{\mu}_i \ge 0, \quad i = 1, \dots, m.
    \]
    Esto, por su vez, es equivalente a la validez de las condiciones de complementaridad
    \[
        g_i(\bar{x}) \le 0, \quad \bar{\mu}_i g_i(\bar{x}) = 0, \quad \bar{\mu}_i \ge 0, \quad i = 1, \dots, m.
    \]
    Notamos, finalmente, que la derivada de la Lagrangiana aumentada en relación a las variables primales se reduce a la derivada de la Lagrangiana clásica:
    \begin{align*}
        0 &= L'_x(\bar{x}, \bar{\lambda}, \bar{\mu}; c) = f'(\bar{x}) + (h'(\bar{x}))^\top\bar{\lambda} + \sum_{i=1}^m \bar{\mu}_i g'_i(\bar{x}) \\
        &= L'_x(\bar{x}, \bar{\lambda}, \bar{\mu}),
    \end{align*}
    lo que concluye la prueba.
\end{proof}

Los dos resultados siguientes son análogos a los de los Lemas 4.5.3 y 4.5.4, respectivamente.

\begin{lemma}\label{v2:lem:4.5.6}
    Sean $f \colon \R^n \to \R$, $h \colon \R^n \to \R^l$ y $g \colon \R^n \to \R^m$ diferenciables en el punto $\bar{x} \in \R^n$.
    Entonces, para todo $c > 0$ y cualesquiera $\lambda \in \R^l$ y $\mu \in \R^m$, $\bar{x}$ es un punto estacionario del problema \eqref{v2:eq:4.211} con función objetivo definida por la fórmula \eqref{v2:eq:4.210}, si, y solamente si, el punto $(\bar{x}, t(\bar{x}))$, donde $t(\bar{x}) \in \R^m$ está definido en \eqref{v2:eq:4.208}, es estacionario para el problema \eqref{v2:eq:4.207} con función objetivo dada por \eqref{v2:eq:4.206}.
\end{lemma}

% [Ejercicio 4.5.12 movido a exercises.tex]

\begin{lemma}\label{v2:lem:4.5.7}
    Para $f \colon \R^n \to \R$, $h \colon \R^n \to \R^l$ y $g \colon \R^n \to \R^m$ cualesquiera, si para algunos $c > 0$, $\lambda \in \R^l$ y $\mu \in \R^m$ el punto $(\tilde{x}, \tilde{t}) \in \R^n \times \R^m$ es una solución local del problema \eqref{v2:eq:4.207} con función objetivo definida por la fórmula \eqref{v2:eq:4.206}, entonces $\tilde{t}$ coincide (ignorando los signos de las coordenadas) con el punto $t(\tilde{x})$, definido de acuerdo con \eqref{v2:eq:4.208}.
\end{lemma}

% [Ejercicio 4.5.13 movido a exercises.tex]

El método de la Lagrangiana aumentada o, aún, el método de los multiplicadores, para el problema \eqref{v2:eq:4.181} es el siguiente.

\noindent \textbf{Algoritmo 4.5.4. (El método de los multiplicadores)} \\
Elegir $\lambda^0 \in \R^l$, $\mu^0 \in \R^m_+$, $c_0 > 0$, y tomar $k := 0$.
\begin{enumerate}
    \item Calcular $x^k \in \R^n$, un punto estacionario del problema \eqref{v2:eq:4.211}, con función objetivo dada por la fórmula \eqref{v2:eq:4.210}, donde $c = c_k, \lambda = \lambda^k$ y $\mu = \mu^k$.
    \item Tomar
    \begin{align}
        \lambda^{k+1} &= \lambda^k + c_k h(x^k), \nonumber \\
        \mu^{k+1}_i &= \max\{0, c_k g_i(x^k) + \mu_i^k\}, \quad i = 1, \dots, m.
        \label{v2:eq:4.214}
    \end{align}
    \item Elegir $c_{k+1} \ge c_k$. \\
    Tomar $k := k + 1$, retornar al Paso 1.
\end{enumerate}

Comentaremos sobre la fórmula \eqref{v2:eq:4.214} para el cálculo de las aproximaciones de los multiplicadores de Lagrange asociados a restricciones de desigualdad. De acuerdo con el esquema del método para el problema \eqref{v2:eq:4.182}-\eqref{v2:eq:4.183} con restricciones transformadas en igualdades (vea el Algoritmo 4.4.3), la fórmula natural sería
\[
    \mu_i^{k+1} = \mu_i^k + c_k(g_i(x^k) + (t_i(x^k, \mu^k, c_k))^2), \quad i = 1, \dots, m.
\]
Más esto es exactamente \eqref{v2:eq:4.214}, si llevamos en cuenta \eqref{v2:eq:4.209}.
El resultado siguiente es la extensión del Teorema 4.4.6 al caso de restricciones mixtas.

\begin{theorem}[Convergencia local del método de los multiplicadores]\label{v2:thm:4.5.7}
    Supongamos que se satisfacen las hipótesis del Teorema 4.5.5.
    Entonces existe una vecindad $U$ del punto $\bar{x}$ y números $\bar{c} \ge 0$, $\delta > 0$ y $M > 0$ tales que:
    \begin{enumerate}[(a)]
        \item Para todo $(\lambda, \mu, c) \in \Delta(\bar{c}, \delta)$, donde
        \[
            \Delta(\bar{c}, \delta) = \left\{ (\lambda, \mu, c) \in \R^l \times \R^m \times \R \;\middle|\; \begin{array}{l} \|\lambda - \bar{\lambda}\| < \delta c, \\ \|\mu - \bar{\mu}\| < \delta c, \\ c > \bar{c} \end{array} \right\},
        \]
        el problema \eqref{v2:eq:4.211}, con función objetivo dada por la fórmula \eqref{v2:eq:4.210}, posee en $U$ un único punto estacionario $x(\lambda, \mu, c)$, y se tiene que
        \[
            \|x(\lambda, \mu, c) - \bar{x}\| \le M\frac{\|\lambda - \bar{\lambda}\| + \|\mu - \bar{\mu}\|}{c},
        \]
        \[
            \|\lambda + c h(x(\lambda, \mu, c)) - \bar{\lambda}\| \le M\frac{\|\lambda - \bar{\lambda}\| + \|\mu - \bar{\mu}\|}{c},
        \]
        \[
            |\max\{0, c g_i(x(\lambda, \mu, c)) + \mu_i\} - \bar{\mu}_i| \le M\frac{\|\lambda - \bar{\lambda}\| + \|\mu - \bar{\mu}\|}{c}
        \]
        para todo $i = 1, \dots, m$;
        \item existe $\bar{\delta} \in (0, \delta]$ tal que para toda sucesión $\{c_k\} \subset \R_+$ no-decreciente y cualquier $\lambda^0 \in \R^l$ y $\mu^0 \in \R^m$ que satisfacen
        \[
            (\lambda^0, \mu^0, c_0) \in \Delta(\bar{c}, \bar{\delta}),
        \]
        las fórmulas
        \[
            \lambda^{k+1} = \lambda^k + c_k h(x(\lambda^k, \mu^k, c_k)),
        \]
        \begin{equation}
            \mu_i^{k+1} = \max\{0, c_k g_i(x(\lambda^k, \mu^k, c_k)) + \mu_i^k\}, \quad i = 1, \dots, m,
            \label{v2:eq:4.215}
        \end{equation}
        donde $k = 0, 1, \dots$, unívocamente definen las sucesiones $\{\lambda^k\}$ y $\{\mu^k\}$, en el sentido de que $\{(\lambda^k, \mu^k, c_k)\} \subset \Delta(\bar{c}, \delta)$, i.e., para todo $k$ el punto $x(\lambda^k, \mu^k, c_k)$ está bien definido e $\{x^k\} \to \bar{x}$, $\{\mu^k\} \to \bar{\mu}$, $\{x(\lambda^k, \mu^k, c_k)\} \to \bar{x}$ $(k \to \infty)$. La tasa de convergencia de las sucesiones $\{\lambda^k\}$ y $\{\mu^k\}$ es lineal, y si $\{c_k\} \to \infty$ $(k \to \infty)$, la tasa es superlineal.
    \end{enumerate}
\end{theorem}

% [Ejercicio 4.5.14 movido a exercises.tex]

Notemos el siguiente hecho, complementario en relación al ítem (b) del Teorema 4.5.7. Como $\{\mu^k\} \to \bar{\mu}$, $\{x(\lambda^k, \mu^k, c_k)\} \to \bar{x}$ $(k \to \infty)$, para todo índice $i \in \{1, \dots, m\} \setminus I(\bar{x})$ y todo $k$ suficientemente grande, se tiene que $c_k g_i(\lambda^k, \mu^k, c_k) + \mu_i^k < 0$. Donde, utilizando también \eqref{v2:eq:4.215}, sigue que $\mu_i^{k+1} = 0$. O sea, las aproximaciones de los multiplicadores que corresponden a restricciones de desigualdad inativas en $\bar{x}$, son exactas (iguales a $\bar{\mu}_i = 0$) después de un número finito de iteraciones.

\subsection{Penalización exacta diferenciable}\label{v2:sec:4.5.5}

Haremos ahora algunos comentarios bien breves sobre penalización exacta diferenciable para el problema con restricciones mixtas \eqref{v2:eq:4.181}. De nuevo, primero introducimos las funciones para el problema \eqref{v2:eq:4.182}-\eqref{v2:eq:4.183}, con restricciones de desigualdad transformadas en restricciones de igualdad. Por ejemplo,
\[
    \tilde{\varphi}(\cdot; c_1, c_2) \colon (\R^n \times \R^m) \times (\R^l \times \R^m) \to \R,
\]
\begin{align}
    \tilde{\varphi}((x, t), (\lambda, \mu); c_1, c_2) &= \tilde{L}((x, t), (\lambda, \mu); c_1) \nonumber \\
    &\quad + \frac{c_2}{2} \|\tilde{L}'_{(x,t)}((x, t), (\lambda, \mu))\|^2 \nonumber \\
    &= L(x, \lambda, \mu) + \sum_{i=1}^m \mu_i t_i^2 \nonumber \\
    &\quad + \frac{c_1}{2} \left( \|h(x)\|^2 + \sum_{i=1}^m (g_i(x) + t_i^2)^2 \right) \nonumber \\
    &\quad + \frac{c_2}{2} \left( \|L'_x(x, \lambda, \mu)\|^2 + 4\sum_{i=1}^m (\mu_i t_i)^2 \right).
    \label{v2:eq:4.216}
\end{align}
Para $c_1, c_2 > 0$, consideraremos el problema
\begin{equation}
    \begin{array}{l}
        \min \tilde{\varphi}((x, t), (\lambda, \mu); c_1, c_2), \\
        ((x, t), (\lambda, \mu)) \in (\R^n \times \R^m) \times (\R^l \times \R^m).
    \end{array}
    \label{v2:eq:4.217}
\end{equation}
No es difícil ver que para todo $x \in \R^n$, $\lambda \in \R^l$ y $\mu \in \R^m$ fijos, el mínimo global en relación a $t \in \R^m$ en \eqref{v2:eq:4.217} es alcanzado en el único (ignorando los signos de las coordenadas) punto $t(x, \mu, c_1, c_2)$, donde
\begin{equation}
    t_i(x, \mu, c_1, c_2) = \begin{cases}
        0, & \text{si } g_i(x) + \mu_i(1 + 2c_2\mu_i)/c_1 > 0, \\
        \sqrt{-(g_i(x) + \mu_i(1 + 2c_2\mu_i)/c_1)}, & \text{caso contrario,}
    \end{cases}
    \label{v2:eq:4.218}
\end{equation}
$i = 1, \dots, m$. Después de algunas transformaciones simples, obtenemos la siguiente función penalizada diferenciable para el problema \eqref{v2:eq:4.181}:
\[
    \varphi(\cdot; c_1, c_2) \colon \R^n \times \R^l \times \R^m \to \R,
\]
\begin{align}
    &\varphi(x, \lambda, \mu; c_1, c_2) \nonumber \\
    &= \min_{t \in \R^m} \tilde{\varphi}((x, t), (\lambda, \mu); c_1, c_2) \nonumber \\
    &= \tilde{\varphi}((x, t(x, \mu, c_1, c_2)), (\lambda, \mu); c_1, c_2) \nonumber \\
    &= f(x) + \langle \lambda, h(x) \rangle + \frac{c_1}{2}\|h(x)\|^2 + \frac{c_2}{2}\|L'_x(x, \lambda, \mu)\|^2 \nonumber \\
    &\quad + \frac{1}{2c_1} \sum_{i=1}^m \left( (\max\{0, c_1 g_i(x) + \mu_i(1 + 2c_2\mu_i)\})^2 \right. \nonumber \\
    &\quad \left. - \mu_i^2(1 + 2c_2\mu_i)^2 - 4c_1 c_2 \mu_i^2 g_i(x) \right).
    \label{v2:eq:4.219}
\end{align}
Consideremos el problema asociado
\begin{equation}
    \min \varphi(x, \lambda, \mu; c_1, c_2), \quad (x, \lambda, \mu) \in \R^n \times \R^l \times \R^m.
    \label{v2:eq:4.220}
\end{equation}

% [Ejercicio 4.5.15 movido a exercises.tex]

El siguiente resultado es análogo al del Lema 4.5.3.

\begin{lemma}\label{v2:lem:4.5.8}
    Sean $f \colon \R^n \to \R$, $h \colon \R^n \to \R^l$ y $g \colon \R^n \to \R^m$ dos veces diferenciables en el punto $\bar{x} \in \R^n$.
    Entonces para todo $c_1 > 0$ y $c_2 \ge 0$, y cualesquiera $\bar{\lambda} \in \R^l$, $\bar{\mu} \in \R^m$, el punto $(\bar{x}, \bar{\lambda}, \bar{\mu})$ es estacionario para el problema \eqref{v2:eq:4.220} con la función objetivo dada por la fórmula \eqref{v2:eq:4.219}, si, y solamente si, el punto $((\bar{x}, \tilde{t}(\bar{x}, \bar{\mu}, c_1, c_2)), (\bar{\lambda}, \bar{\mu}))$, donde $t(\bar{x}, \bar{\mu}, c_1, c_2) \in \R^m$ está definido en \eqref{v2:eq:4.218}, es estacionario para el problema \eqref{v2:eq:4.217} con función objetivo definida por \eqref{v2:eq:4.216}.
\end{lemma}

% [Ejercicio 4.5.16 movido a exercises.tex]

Finalmente, el resultado siguiente generaliza la afirmación (b) del Teorema 4.4.7.

\begin{theorem}\label{v2:thm:4.5.8}
    Supongamos que se satisfacen las hipótesis del Teorema 4.5.5.
    Entonces existe un número $\bar{c}_2 > 0$ tal que, si $c_2 \in (0, \bar{c}_2]$, entonces existe una vecindad $U$ del punto $(\bar{x}, \bar{\lambda}, \bar{\mu})$ y un número $\bar{c}_1(c_2) \ge 0$ pueden ser escogidos de tal manera que para todo $c_1 > \bar{c}_1(c_2)$ el problema \eqref{v2:eq:4.220}, con función objetivo definida por la fórmula \eqref{v2:eq:4.219}, posee en $U$ un único punto estacionario, que es $(\bar{x}, \bar{\lambda}, \bar{\mu})$.
\end{theorem}

% [Ejercicio 4.5.17 movido a exercises.tex]